\documentclass[11pt,reqno]{amsart}

\usepackage[utf8]{inputenc}
\usepackage[T1]{fontenc}
\usepackage{lmodern}
\usepackage{amsmath,amssymb,amsthm,booktabs,hyperref}
\usepackage[margin=1in]{geometry}

\newtheorem{theorem}{Theorem}[section]
\newtheorem{proposition}[theorem]{Proposition}
\newtheorem{lemma}[theorem]{Lemma}
\newtheorem{corollary}[theorem]{Corollary}
\theoremstyle{remark}
\newtheorem{remark}[theorem]{Remark}

\newcommand{\ind}{\mathbf{1}}
\newcommand{\normF}{\widetilde{F}}
\newcommand{\normf}{\widetilde{f}}
\newcommand{\normD}{\widetilde{D}}
\newcommand{\cA}{\mathcal{A}}
\newcommand{\cT}{\mathcal{T}}

\numberwithin{equation}{section}

\title[On Proposition $(1-1.67539)$ for the Twin Prime Problem]{A Self-Contained Proof of Proposition $(1-1.67539)$\\ for the Twin Prime Problem}

\author{Wei Lindong}
\address{Independent Researcher}
\email{744140495@qq.com}
\hypersetup{
  pdftitle={A Self-Contained Proof of Proposition (1-1.67539) for the Twin Prime Problem},
  pdfauthor={Wei Lindong},
  pdfsubject={Analytic number theory; sieve methods},
  pdfkeywords={twin prime problem, Chen's theorem, weighted sieve, level of distribution, integer certificate}
}

\date{November 26, 2026}

\subjclass[2020]{Primary 11N36; Secondary 11N05, 11P32}
\keywords{Twin prime problem, Chen's theorem, weighted sieve, level of distribution, directed integer certificate}

\begin{document}

\begin{abstract}
For $1 \le a \le 2$, Proposition $(1-a)$ asserts that there exist infinitely many primes $p$ such that $p+2 = rq$ with $q$ prime, $r=1$ or prime, and $r \le q^{a-1}$. Li and Liu recently established Proposition $(1-1.75)$ unconditionally. By combining Runbo Li's published lower bound $\pi_{1,2}(x) \ge 1.2759\,C_2 x/(\log x)^2$ with a self-contained positive upper-sieve majorant at distribution level $5/8$ (derived from Pascadi's fixed-residue convolution estimate) and a directed pure-integer interval certificate at denominator $2^{80}$, we prove Proposition $(1-1.67539)$ unconditionally. More precisely, writing $A_b(x)$ for the number of primes $p \le x$ satisfying $p+2 = rq$ with $r \le q^b$, we obtain $A_{67539/100000}(x) > 0.0000466\,C_2 x/(\log x)^2$ for all sufficiently large $x$. We also record an elementary one-page two-factor shortcut yielding Proposition $(1-1.7078)$ and a rank-three baseline yielding Proposition $(1-1.6767)$.
\end{abstract}

\maketitle

\section{Introduction and main results}\label{sec:intro}

Let
\[
C_2 := 2\prod_{\ell > 2}\left(1 - \frac{1}{(\ell-1)^2}\right) \approx 1.3203236317,
\qquad
H(x) := \frac{C_2 x}{(\log x)^2},
\]
where $\ell$ runs over primes. For $0 < b \le 1$, let $A_b(x)$ count the number of primes $p \le x$ admitting a representation
\[
p + 2 = rq,
\]
where $q$ is prime, $r$ is either $1$ or prime, and $r \le q^b$. In the notation of Li and Liu~\cite{LiLiu2026}, Proposition $(1-a)$ ($1 \le a \le 2$) states that $A_{a-1}(x) \to \infty$ as $x \to \infty$; thus $a=2$ corresponds to Chen's theorem~\cite{Chen1973} and $a=1$ corresponds to the twin prime conjecture.

In~\cite[Theorem~1.2]{LiLiu2026}, Li and Liu proved Proposition $(1-1.75)$ unconditionally, obtaining
\[
A_{3/4}(x) > 0.036\,H(x)
\]
for sufficiently large $x$. Under a weighted Elliott--Halberstam conjecture $\mathrm{WEH}(0.999)$, they improved the exponent to $(1-1.4)$~\cite[Theorem~1.4]{LiLiu2026}. Our main result strengthens the unconditional bound from $1.75$ to $1.67539$.

\begin{theorem}\label{thm:main}
Proposition $(1-1.67539)$ holds unconditionally. Specifically, for $b = 67539/100000$,
\begin{equation}\label{eq:main_margin}
\liminf_{x\to\infty}\frac{A_b(x)}{H(x)}
\ \ge\ \frac{56380964408304863785}{1208925819614629174706176}
\ >\ 0.0000466.
\end{equation}
In particular, $A_{67539/100000}(x) > 0.0000466\,H(x)$ for every sufficiently large $x$. Without the fourth-rank cap replacement, the rank-three construction alone establishes Proposition $(1-1.6767)$, while an elementary two-factor shortcut proves Proposition $(1-1.7078)$.
\end{theorem}

\subsection{Comparison with related work}
It is helpful to distinguish Theorem~\ref{thm:main} from related results on shifted primes $p+2$:
\begin{itemize}
\item \emph{Chen's theorem density.} Runbo Li~\cite[Theorem~1.3]{RunboLi2026} established the lower bound $\pi_{1,2}(x) \ge 1.2759\,H(x)$ for $\pi_{1,2}(x) = \#\{p \le x : p+2 \in P_2\}$, using Pascadi's distribution results~\cite{Pascadi2025}. Our proof uses Li's $1.2759$ lower bound directly as a modular input.
\item \emph{Largest prime factor.} Baker and Harman~\cite{BakerHarman1996} (discussed in~\cite[\S 1.4]{LiLiu2026}) showed that $P^+(p+2) > p^{0.676}$ infinitely often, giving $p+2 = rq$ with $q$ prime and $r \le p^{0.324} \le q^{0.48}$. However, the cofactor $r$ in Baker--Harman's theorem is an arbitrary integer that may have three or more prime factors, whereas Theorem~\ref{thm:main} retains Chen's condition $\Omega(p+2) \le 2$ ($r=1$ or $r$ prime).
\item \emph{Twin-prime upper bounds and Goldbach analogues.} Building on Lichtman~\cite{Lichtman2023} (who obtained $3.229\,H(x)$ at level $66/107$), Pascadi~\cite[Corollary~1.4]{Pascadi2025} proved $\pi_2(x) \le (3.203 + o(1))H(x)$ using an unconditional $5/8$ distribution level~\cite[Proposition~4.4 and Theorem~5.4]{Pascadi2025}. We adapt Chan's positive prefix sieve and subtraction of excluded semiprimes from the Chen count~\cite[\S\S 4 and 8 and Appendix~F]{Chan2026}. At the fixed residue $2$, we use Pascadi's $5/8$ estimate, verify the modified rank-three and rank-four modulus factorizations for prime pairs $rq$, and certify the resulting density bounds with directed integer arithmetic.
\end{itemize}

\subsection{Inputs and organization}
The external analytic inputs are Runbo Li's Theorem~1.3~\cite{RunboLi2026}; Pascadi's fixed nonzero residue convolution estimate~\cite[Proposition~4.4]{Pascadi2025} (which itself relies on Deshouillers--Iwaniec~\cite{DeshouillersIwaniec1982}, Watt, and Pascadi's exceptional Maass-form large sieve); Iwaniec's well-factorable linear Rosser sieve and the fundamental lemma~\cite{Iwaniec1980,Iwaniec1980b,FriedlanderIwaniec2010}; and the classical prime number theorem, Mertens's estimates, prime and M\"obius Siegel--Walfisz bounds, the multiplicative large sieve, and the Weil bound for Kloosterman sums. No external distribution theorem or density bound for the constructed coefficient is assumed.

In the remainder of Section~\ref{sec:intro}, we reduce $A_b(x)$ to subtracting the bad semiprimes $B_b(x)$ and present the one-page shortcut to $(1-1.7078)$. Section~\ref{sec:tree} defines the rank-three and rank-four upper-sieve trees and proves the pointwise first-exit identity. Section~\ref{sec:packing} verifies the capacity conditions for Pascadi~\cite[eq.~(4.15)]{Pascadi2025} (including the $C<1$ two-factor repair and three-slot $W$ packing) and realizes a single finite signed coefficient $\rho_x$. Section~\ref{sec:distribution} establishes the uniform prime-pair distribution estimate for $\rho_x$ at fixed residue $2$. Sections~\ref{sec:density} and~\ref{sec:bad_pairs} compute the limiting density, bound the small-prime tail, and transfer the majorant to $B_b(x)$. Section~\ref{sec:certificate} describes the standalone integer certificate script \texttt{verify\_twin\_upper\_sieve\_cert.py} and records the certified rational bounds.

\subsection{The subtraction identity and an elementary shortcut}
Let $B_b(x)$ count the primes $p \le x$ such that $p+2 = rq$ with $r,q$ prime and $q^b < r \le q$. Ordering the two prime factors ensures that every composite $P_2$ remainder has a unique factorization $r \le q$ (squares contribute to $B_b$), while prime remainders ($r=1$) belong to $A_b$. Allowing for at most one exception at $p=2$, Runbo Li's lower bound~\cite[Theorem~1.3]{RunboLi2026} gives
\begin{equation}\label{eq:sub}
\liminf_{x\to\infty}\frac{A_b(x)}{H(x)}
\ \ge\ \frac{12759}{10000} - \limsup_{x\to\infty}\frac{B_b(x)}{H(x)}.
\end{equation}

In odd moduli $d$, the congruence $rq-2 \equiv 0 \pmod d$ reads $rq \equiv 2 \pmod d$. The two-prime sequence has local density $g(\ell) = 1/(\ell-1)$ for odd primes $\ell$, with no factor at $2$. Since $(1-g(\ell))/(1-1/\ell) = 1 - (\ell-1)^{-2}$, putting $P(y) = \prod_{2 < \ell < y}\ell$ and $V(y) = \prod_{2 < \ell < y}(1-g(\ell))$ yields
\begin{equation}\label{eq:mertens_V}
V(x^t) \sim \frac{e^{-\gamma}C_2}{t\log x}
\end{equation}
uniformly for $t$ in any positive compact interval. The factor $2$ from omitting $\ell=2$ is already contained in $C_2$, matching~\cite[(6.6)--(6.10)]{LiLiu2026} and~\cite[Corollary~1.4]{Pascadi2025}.

\begin{proposition}[Two-factor shortcut]\label{prop:shortcut}
Proposition $(1-1.7078)$ holds unconditionally; more precisely, for $b = 3539/5000$,
\[
\liminf_{x\to\infty}\frac{A_b(x)}{H(x)}
\ \ge\ \frac{377883443722357593132}{1208925819614629174706176}
\ >\ 0.000312.
\]
\end{proposition}

\begin{proof}
We apply Li and Liu's fixed-residue bilinear distribution lemma~\cite[Lemma~3.4]{LiLiu2026} (originating in Bombieri--Friedlander--Iwaniec and Wu) to one-bounded well-factorable modulus weights. Take $N$ in that lemma to be the shorter dyadic scale (corresponding to $r$); prime indicators satisfy its Siegel--Walfisz and roughness hypotheses. In the region $q^b < r \le q$ with $rq \asymp X$ and $b > 2/3$, we have $\nu := \log N/\log X \in [2/5, 1/2]$ for large $x$, allowing bounded dyadic endpoint errors and choosing the smaller scale if $r$ exceeds $q$ near the diagonal. The table in~\cite[Lemma~3.4]{LiLiu2026} gives distribution exponent $1 - \nu - \epsilon$. On each dyadic chart we apply the ordinary one-bounded well-factorable upper linear sieve (Section~\ref{sec:packing} below shows how to realize this greedily atom-by-atom via a single constant cubic cap below $1-\nu-\epsilon$, reserving the small-prime attachment $x^{\sqrt{\eta}}$ in the larger group). With $t = \log r/\log x$ and sifting threshold at a third of the level exponent,~\eqref{eq:mertens_V} and $F(3) = 2e^\gamma/3$ yield limiting upper density $2/(1-t)$ in units of $C_2/\log x$ after letting fixed reserves tend to zero. Smooth separation as in Section~\ref{sec:distribution} preserves prime Siegel--Walfisz by partial summation. The region $rq < \sigma x$ contributes $O(\sigma H(x))$ using a fixed level $1/2$ after discarding $rq < x/(\log x)^5$. Integrating against $dt/(t(1-t))$ gives
\[
\limsup_{x\to\infty}\frac{B_b(x)}{H(x)}
\ \le\ \int_{b/(1+b)}^{1/2}\frac{2\,dt}{t(1-t)^2}
\ =\ 2(1 - b - \log b),
\]
since an antiderivative of $1/(t(1-t)^2)$ is $\log(t/(1-t)) + (1-t)^{-1}$. Substituting into~\eqref{eq:sub} and evaluating at $b = 3539/5000$ via Section~\ref{sec:certificate} completes the proof.
\end{proof}

\section{The positive upper-sieve tree and cap rules}\label{sec:tree}

Write $L = \log x$, $\rho = 1/5$, $T = 5/8$, and
\[
v(u) := \min\left\{\frac{1+u}{2},\, 1 - \frac{3u}{2}\right\}.
\]
Order the prime factors of any odd squarefree modulus decreasingly as $\ell_1 > \ell_2 > \dots > \ell_r > 2$, and set $t_i = \log \ell_i/L$ and $S_j = \sum_{i \le j} t_i$ ($S_0 = 0$). Fix a root cutoff $\rho' < \rho$ arbitrarily close to $\rho$. An odd-rank prefix $(\ell_1,\dots,\ell_j)$ is retained precisely when its local cubic test satisfies $S_{j-1} + 3t_j < \lambda$, where $\lambda$ is the current cap; every even child of a retained odd prefix is retained automatically.

We start at rank one with cap $T$. At rank three, writing $(t_1,t_2,t_3)$ before testing the third cube, replace the cap by the ideal third-rank cap $J(t_1,t_2,t_3)$, defined as the maximum of:
\begin{itemize}
\item $v(t_1)$, $v(t_1+t_2)$, and $v(t_1+t_2+t_3)$;
\item $v(t_1+t_3)$, if $v(t_1+t_3) > t_1 + t_3 + 2t_2$;
\item $v(t_2+t_3)$, if $v(t_2+t_3) > t_2 + t_3 + 2t_1$;
\item $\min\{5/8 - 3t_3/8,\, 1 - 2t_2\}$.
\end{itemize}
Notice the explicit $1 - 2t_2$ guard in the last branch, which is essential here since $2\rho = 2/5 > 1 - 5/8$. No unguarded nonleading complementary square is used.

At an admitted rank-four prefix $t = (t_1,t_2,t_3,t_4)$ with $t_4 \ge 3/100$, examine every subset $I \subseteq \{1,2,3,4\}$ containing $4$ and at least one other index, and put $u_I := \sum_{i \in I} t_i$. Raise the inherited cap to the largest $M = v(u_I) > J$ over those candidates satisfying
\begin{equation}\label{eq:rank4_guard}
u_I + \sum_{\substack{i < j\\ i \notin I}} t_i + 2t_j < M \qquad (j \le 4,\ j \notin I).
\end{equation}
If no candidate satisfies~\eqref{eq:rank4_guard}, retain $J$. Since rank four is even, this update occurs without a retention test at rank four. All subsequent odd ranks use the inherited cap. In Section~\ref{sec:packing}, strictly lowered piecewise-constant versions of these ideal caps will be frozen for the arithmetic construction.

\begin{lemma}[First-exit identity]\label{lem:first_exit}
For any finite prime set and any multiplicative function $g$ satisfying $0 \le g(\ell) \le 1$, the retained prefix tree $\cT$ satisfies
\begin{equation}\label{eq:first_exit}
\sum_{d \in \cT}\mu(d)g(d)
= \prod_\ell (1 - g(\ell))
+ \sum_{e\ \mathrm{first\ omitted\ prefix}} g(e)\prod_{\ell < \ell_{\mathrm{last}}(e)}(1 - g(\ell)).
\end{equation}
In particular, taking $g(\ell) = \ind_{\ell \mid n}$ shows that $\sum_{d \mid n,\, d \in \cT}\mu(d) \ge \ind_{(n,\prod\ell)=1} \ge 0$ for every integer $n$.
\end{lemma}

\begin{proof}
Every divisor $d \notin \cT$ has a unique first omitted prefix $e$, which necessarily has odd rank so that $\mu(e) = -1$. Summing $\mu(d)g(d)$ over all squarefree continuations of $e$ by primes $\ell < \ell_{\mathrm{last}}(e)$ yields $-g(e)\prod_{\ell < \ell_{\mathrm{last}}(e)}(1-g(\ell))$. Subtracting these disjoint sums from $\prod_\ell(1-g(\ell))$ gives~\eqref{eq:first_exit}. Raising a cap after an even prefix or selecting it before its own odd test preserves the odd parity of every first exit.
\end{proof}

\section{Capacity packing and finite realization}\label{sec:packing}

For a slot index $j$, the \emph{full square} is $S_{j-1} + 2t_j$. For a selected index subset $I$, the \emph{complementary squares} are $u_I + \sum_{i<j,\, i\notin I} t_i + 2t_j$ for $j \notin I$, as in~\eqref{eq:rank4_guard}; after the largest index of $I$, each complementary square coincides with the full square. Any passing odd cube $S_{j-1} + 3t_j < \lambda$ implies both its own full square and the next full square $S_j + 2t_{j+1} \le S_{j-1} + 3t_j < \lambda$.

For the initial checks of $J$: the first two full squares are bounded by $3t_1$, while $v(t_1) = (1+t_1)/2 > 3t_1$ holds strictly for all $t_1 < 1/5$. Moreover, a passing third cube controls both the third and fourth full squares even if the rank-four cap is subsequently raised. For $I=\{1\}$, complementary squares reduce to full squares; for $I=\{1,2\}$, the only complementary square up to rank three is a full square; for $I=\{1,2,3\}$, there is none; for $I=\{1,3\}$ and $I=\{2,3\}$, the sole extra square is explicitly guarded in the definition of $J$; and for rank four,~\eqref{eq:rank4_guard} enforces every extra square directly. These inequalities persist under appending smaller primes.

\begin{lemma}[Greedy modulus packing]\label{lem:packing}
Let $X$ be the complete product scale, and let $U$ be a selected medium-variable scale satisfying $X^{2\zeta} \le U \le X^{1/3+\zeta/2}$, with $0 < \zeta < 10^{-5}$ sufficiently small. Convert exponents to base $X$, let $\ell>1/2$ be a covering exponent, and put $D = X^{\ell - 2\zeta}$ and $A = U X^{-\zeta}$. Fix a covering reserve $\delta>4\zeta$. Require all full square exponents of the box atom to be at most $\ell-\delta$, and impose the same bound on complementary square exponents in case~(a). Assume one of the two conditions below holds.
\begin{enumerate}
\item[(a)] $\ell < v(\beta)$ for a selected subset $I$ of upper endpoint $B = X^\beta$; or
\item[(b)] $(1+t_1)/2<\ell<5/8-3t_3/8$ and $2t_2\le1-\ell-\delta$.
\end{enumerate}
Then the atom admits a partition into three modulus groups with upper endpoints $Q_1, Q_2, Q_3 \ge 1$ satisfying Pascadi's system~\cite[eq.~(4.15)]{Pascadi2025}:
\begin{equation}\label{eq:pascadi_cond}
Q_1 \le U X^{-\epsilon}, \qquad U^2 Q_2 Q_3^2 \le X^{1 - 15\epsilon}, \qquad U^2 Q_2^5 Q_3^2 \le X^{2 - 40\epsilon},
\end{equation}
provided $\epsilon > 0$ is sufficiently small in terms of $\zeta$ and the covering margins.
\end{lemma}

\begin{proof}
In case~(a), the condition $\ell < v(\beta)$ is equivalent to
\begin{equation}\label{eq:beta_interval}
b_- := 2\ell - 1 < \beta < \frac{2 - 2\ell}{3} =: b_+.
\end{equation}
If $C := D/(AB) \ge 1$, place the selected slots $I$ in group two ($Q_2 = B$) and pack the remaining slots greedily in order into capacities $A$ and $C$. If some slot of exponent $s$ fits neither, the product of the remaining capacities is $< X^{2s}$, contradicting the complementary square at that slot. With $(Q_1,Q_2,Q_3) = (A,B,C)$, replacing $U$ by $A X^\zeta$ bounds the two monomials in~\eqref{eq:pascadi_cond} by $X^{2\zeta} D^2/B$ and $X^{2\zeta} D^2 B^3$. By~\eqref{eq:beta_interval}, $D^2/B \le X^{2\ell - 4\zeta - (2\ell-1)} = X^{1-4\zeta}$ and $D^2 B^3 \le X^{2\ell - 4\zeta + 2 - 2\ell} = X^{2-4\zeta}$, leaving fixed power savings $2\zeta$ in both bounds.

If instead $C = D/(AB) < 1$, repack \emph{all} slots greedily into $(A, D/A)$ using the full squares and leave group three empty ($Q_3 = 1$). Note that $A \ge X^\zeta \ge 1$ and $D/A \ge X^{\ell - 1/3 - 3\zeta/2} \ge 1$ since $\ell > 1/2$. Because $Q_2 = D/A < B$, the second monomial $X^{2\zeta} D^2 Q_2^3$ strictly decreases from $X^{2\zeta} D^2 B^3$, while the first becomes $U^2 (D/A) = X^{2\zeta} A D \le X^{1/3 + \ell - \zeta/2}$. Both cases give $\ell\le 5/8$, so the latter exponent is at most $23/24-\zeta/2<1-4\zeta$ for the chosen $\zeta$. (This is the same two-factor repair as~\cite[Lemma~5.6]{Pascadi2025}.)

In case~(b), initialize capacities $(A, X^{b_-}, D/(AX^{b_-}))$. If the third capacity is below $1$, apply the same two-group repair into $(A, D/A)$. Otherwise, place the largest prime in group two (since $t_1 < 2\ell - 1 = b_-$). The second prime fits into group one or three because its square is at most $X^{1-\ell-\delta}<X^{1-\ell-2\zeta}$, the product of their capacities. Continue inserting slots greedily in order. If all fit, the estimates for case~(a) apply with $\beta=b_-$, since the nonstrict bounds $b_-\le\beta\le b_+$ still leave the explicit $2\zeta$ savings. At the first slot $\ell_j$ of exponent $s \le t_3$ that fits no group, adjoin $\ell_j$ to group two, obtaining a new product $B'$ with exponent in $[b_-, b_- + t_3]$. Since $b_+ - b_- = (5 - 8\ell)/3 > t_3$,~\eqref{eq:beta_interval} still holds for $B'$. Moreover, if $(a_1,a_2,a_3)$ were the partial products before $\ell_j$, then $A < a_1 X^s$ and the full square gives $a_1 a_2 a_3 X^{2s} \le D$, whence $A B' a_3 < D$ for $B' = a_2 X^s$. Thus $a_3$ fits inside $D/(A B')$, and freezing group two lets the remaining slots be packed into groups one and three via the full squares. Both monomial bounds in~\eqref{eq:pascadi_cond} follow as before.
\end{proof}

\begin{remark}
In the finite realization below, small-prime weights up to $x^{\sqrt{\eta}}$ are attached to group three; choosing $\eta$ sufficiently small relative to the covering reserves absorbs an extra $x^{\sqrt{\eta}}$ factor in $Q_3$ into the strict margins of~\eqref{eq:pascadi_cond}.
\end{remark}

\subsection{Construction of a single finite coefficient}
Fix $a_*>0$, the density allowance $u_0>0$, the root cutoff $\rho'<\rho$, and the third-slot tail cutoff $\delta_0>0$. Work first in base $x$. On each of finitely many coarse cells freeze the branch and its physical cap $\lambda_{\rm phys}$. Where packing is used, choose a covering exponent $\ell>\lambda_{\rm phys}$ below the relevant ideal endpoint. The positive gaps arise as follows.
\begin{itemize}
\item \emph{Small $t_1$ and ranks at most two.} Use a fixed cap below $1/2$ for small $t_1$. A rough prefix of rank at most two has total exponent at most $2\rho<1/2$. Both gaps allow the absolute mean value bound.
\item \emph{Fallback.} At the third-slot tail and trimmed walls use $I=\{1\}$ and a slightly lowered $v(t_1)$. Away from small $t_1$, $v(t_1)>1/2$ uniformly, and $v(t_1)-3t_1\ge(1-5\rho')/2>0$ makes room for the first two full squares below the physical cap. Test the third cube on the interval upper endpoints as $t_1^++t_2^++3t_3^+<\lambda_{\rm phys}$. This controls full squares three and four, and later full squares follow from the later odd cubes. For $I=\{1\}$ the complementary squares are full squares.
\item \emph{Compact third-rank cells.} Keep a selected set only with a fixed gap in its extra guards. Lower its physical cap by less than these gaps. Trim near ties so that the physical cap still exceeds the fallback cap. The first two full squares are controlled by the fallback bound and the other full squares by passing odd cubes. On a $W$ cell choose $\ell$ above both $\lambda_{\rm phys}$ and $(1+t_1)/2$ and below both $5/8-3t_3/8$ and $1-2t_2$ throughout the cell. The second upper bound gives a separate gap in $2t_2<1-\ell$.
\item \emph{Compact fourth-rank cells.} Trim near equality in~\eqref{eq:rank4_guard} and in the increase over the inherited cap. The new cap may be lowered slightly while preserving these gaps. The inherited passing third cube controls full squares three and four. The new cap controls subsequent odd cubes. The extra complementary squares through slot four follow from~\eqref{eq:rank4_guard}; beyond the last selected slot they are full squares. Elsewhere retain the parent rule.
\end{itemize}
The gap $\ell-\lambda_{\rm phys}$ makes each square bound by the physical cap strict relative to $\ell$. On the finitely many retained cells there is a common $\delta_1>0$ smaller than these covering gaps, the guard and analytic endpoint gaps, and the gaps for the absolute mean value bound. Trimming switching walls and using the fallback at small $t_3$ have arbitrarily small density cost by Section~\ref{sec:density}. The approximation retains a positive gap above $t_1+t_2$, since $J-t_1-t_2\ge(1-3\rho)/2$ before lowering.

Put $\delta=\delta_1/4$, choose $0<\zeta<\min(10^{-5},\delta/4,a_*/24)$, then take $\eta,\epsilon>0$ small enough for the packing and attachment. At fixed $\eta$ the rough tree on primes at least $x^\eta$ has bounded rank. Refine the common interval alphabet so that upper endpoints use at most $\delta_1/4$ of the initial covering and guard gaps. Test odd cubes at the physical caps on those endpoints and retain even children automatically. Finally choose $\kappa$ as in Proposition~\ref{prop:distribution}. On charts $X\gg x^{1-\kappa}$, the factor $\log x/\log X=1+O(\kappa+L^{-1})$ uses at most another $\delta_1/4$ of these gaps. For later full squares, rescale the bound by the physical cap directly, using that the cap is bounded. The resulting base-$X$ atoms satisfy the reserve $\delta$ of Lemma~\ref{lem:packing}, including $2t_2\le1-\ell-\delta$ in case~(b).

For each rough atom of sign $+$ or $-$, attach the ordinary upper ($\lambda^+$) or lower ($\lambda^-$) small-prime weights on $\ell < x^\eta$ of support at most $x^{\sqrt{\eta}}$ and $|\lambda^\pm| \le 1$. By the fundamental lemma, small-prime density equals $V(x^\eta)(1 + O(e^{-c/\sqrt{\eta}}))$ uniformly in $x$, while the absolute rough density is $O(1/\eta)$, giving attachment error $O(e^{-c/\sqrt{\eta}}\eta^{-2} C_2/L)$. Moreover, since
\[
\lambda^-(n) \le \ind_{(n, P(x^\eta))=1} \le \lambda^+(n) \qquad (\forall n \in \mathbb{Z}),
\]
multiplying by the atom's sign and rough indicator produces a pointwise majorant of the rough-tree sum times $\ind_{(n,P(x^\eta))=1}$, hence majorizes $\ind_{(n,P(x^{\rho'}))=1}$ by Lemma~\ref{lem:first_exit}. Unique prime factorization across $x^\eta$ ensures $|\rho_x(d)| \le 1$, supported on odd squarefree $d \le x^{5/8+\sqrt{\eta}}$. Rough pieces of rank $\le 2$ have support $\le x^{2\rho+\sqrt{\eta}}$. In the small-$t_1$ branch, the support exponent is bounded by its fixed cap plus $\sqrt{\eta}$. Both stay below $1/2-\epsilon_0$ for sufficiently small $\eta,\epsilon_0$.

By Lemma~\ref{lem:first_exit}, multiplying the unrounded rough-tree density by $V(x^\eta)$ isolates the subset of full exits whose last prime is $\ge x^\eta$, which is at most the full unrounded density. For fixed $\eta$, refining the common alphabet reduces the discrepancy to $O_{\mathrm{mesh}\to 0}(C_2/L)$ via Mertens's harmonic measure $dt/t$ on finitely many compact cells.

Finally, for each rough atom with interval multiplicities $m_I$, its ordered squarefree coefficient equals $\ind_{\mathrm{sf}}\,\mathcal{N}/\prod_I m_I!$, where $\mathcal{N}$ is the unrestricted interval convolution. Dividing each of the three packed groups by its own $\prod_I m_{I,j}!$ makes every group coefficient $1$-bounded even on repeated primes, leaving a fixed external scalar. Attaching $\lambda^\pm$ to group three preserves $1$-boundedness by disjoint prime universes. Removing the global squarefree restriction only during analytic comparison introduces a progression error at most
\begin{equation}\label{eq:sf_error}
x^{o(1)} \sum_{\substack{d \le x^{5/8+\sqrt{\eta}}\\ \exists \ell \ge x^\eta:\, \ell^2 \mid d}} \left(\frac{x}{d} + 1 + \frac{x}{\varphi(d)}\right)
\ \ll\ x^{1 - \eta + o(1)},
\end{equation}
which handles $rq - 2 = 0$ as well (no divisor bound at zero is needed).

\section{Distribution on prime pairs}\label{sec:distribution}

Set $\Delta_d(n) := \ind_{n \equiv 2 \pmod d} - \ind_{(n,d)=1}/\varphi(d)$.

\begin{proposition}\label{prop:distribution}
Fix $a_*>0$ first and construct $\rho_x$ as in Section~\ref{sec:packing}, with reserves chosen for that $a_*$. For any smooth weight $w(r,q)$ supported on $r,q \gg x^{a_*}$ and $rq = O(x)$ whose mixed Euler derivatives are bounded by fixed powers of $L = \log x$, we have
\begin{equation}\label{eq:distrib}
\sum_d \rho_x(d) \sum_{r,q\ \mathrm{prime}} w(r,q)\Delta_d(rq)
\ =\ O_A(x L^{-A}) \qquad (\forall A > 0).
\end{equation}
\end{proposition}

\begin{proof}
Fix the maximal divisor-function order $K_{\rm mom}$ used in the moment estimates before choosing $\kappa$. The orders are fixed by the number of Heath--Brown slots, the modulus factors and the squared coefficient bound in truncation. The requested saving $A$ changes logarithmic cutoffs and the number of integrations by parts, so $K_{\rm mom}$ is independent of $A$. Choose $0<\kappa<1/(2K_{\rm mom})$ small enough also to preserve the covering reserves and give total support $D_0\le X^{5/8+\xi}$ on each chart, with fixed small $\xi>0$. Smoothly discarding $rq<2x^{1-\kappa}$ costs $x^{1-\kappa+o(1)}$ by $\tau(rq-2)$ and $\sum_{d\ll x}1/\varphi(d)\ll L$.

On an interval of length $O(X)$ inside $[1,O(x)]$ the $d_m$ sum is $O_m((X+x^{1-1/m})L^{m-1})$. Choose a largest factor in each ordered $m$-tuple. The product $d$ of the other factors satisfies $d\ll_m x^{1-1/m}$, and the last factor has $O(X/d+1)$ choices. For $m\le K_{\rm mom}$, $x^{1-1/m}\ll X$ by the choice of $\kappa$ and $X\gg x^{1-\kappa}$. The same bound applies after the shift by $2$. Cauchy--Schwarz and $\tau^h\le d_{2^h}$ give, for the needed $j,k$,
\begin{equation}\label{eq:shifted_div}
\sum_{n \asymp X} \tau(n)^j \tau(|n-2|)^k \ll X L^{O_{j,k}(1)},
\end{equation}
while coprime means add at most $L^{O_J(1)}$ under $\tau(d)^J/\varphi(d)$.

For pieces with support $D_0 \le X^{1/2-\epsilon_0}$, separate the smooth weight in logarithmic coordinates and apply the multiplicative large sieve: writing $d = gh$ with primitive conductor $g > 1$ and retaining $(n,h)=1$ on both prime factors, Cauchy--Schwarz bounds the $G < g \le 2G$ contribution with weight $1/\varphi(g)$ by $L^{O(1)}(X/G + \sqrt{X}(\sqrt{R}+\sqrt{U}) + G\sqrt{X})$. Summing $G > L^B$ over $h \le D_0/g$ via $1/\varphi(gh) \le 1/(\varphi(g)\varphi(h))$ and $\sum_h 1/(h\varphi(h)) < \infty$, and handling $g \le L^B$ by prime Siegel--Walfisz (removing $\ell \mid h$ at cost $O(\omega(h)L^{O(1)})$), gives $O_A(XL^{-A})$.

For the high-support pieces, fix a common roughness cutoff $z_* = (C X_0)^{1/(\log\log(C X_0))^3}$ across all subcharts of each original dyadic chart of scale $X_0$; for large $X_0$ this exceeds $X^{1/(\log\log X)^3}$. Replace each prime indicator by $\Lambda(m)/\log m$ with $P^-(m) > z_*$ (proper prime powers $\gg x^{a_*}$ cost $O(x^{-a_*/2+o(1)})$). At each original scale $R_i$, put $Z_i = (2R_i)^{1/3}$ and $m_i = \mu\ind_{[1,Z_i]}$; expanding $(\delta - 1*m_i)^{*3}*\Lambda = 0$ below $2R_i$ yields the order-three Heath--Brown identity
\begin{equation}\label{eq:hb3}
\Lambda = \sum_{j=1}^3 (-1)^{j-1}\binom{3}{j} m_i^{*j} * 1^{*(j-1)} * \log \qquad (n \le 2R_i).
\end{equation}
Factor roughness over the at most six variables per prime, dyadically subdivide each slot, and separate the smooth profile and $\log$ weights by Fourier inversion in logarithmic coordinates with two integrations by parts per variable (arbitrarily many to truncate frequencies at $L^B$, controlling tails via~\eqref{eq:shifted_div}).

\smallskip
\noindent\emph{Case 1: A medium slot exists.} Suppose some dyadic slot has length $U \in [X^{2\zeta}, X^{1/3+\zeta/2}]$. Take it for Pascadi's short sequence $\alpha_n$ (a bounded rough $1$ or rough $\mu$ times a logarithmic twist on an interval). To verify Siegel--Walfisz~\cite[Definition~4.3]{Pascadi2025}:
\begin{itemize}
\item For rough $1$, use the even and odd Brun polynomials $U_z, L_z$ below, expanding $d \le y_* = X^{o(1)}$ and counting multiples in each class modulo polylogarithmic $q$; the two bounds match up to $O_A(U L^{-A}) + O(y_* L^{O(1)})$.
\item For rough $\mu$, expand $\mu(n)\ind_{P^-(n)>z_*} = \sum_{d \mid n,\, P^+(d)\le z_*} \mu(n/d)$. Terms with $d > U^{1/2}$ contribute at most $U \exp(-\log U/(2\log z_*))\prod_{\ell \le z_*}(1-\ell^{-1+1/\log z_*})^{-1} \ll_A U L^{-A}$ by Rankin's trick and Mertens; terms with $d \le U^{1/2}$ satisfy Möbius Siegel--Walfisz at length $U/d \ge U^{1/2}$.
\end{itemize}
Auxiliary coprimality $(n,h)=1$ deletes at most $O((U/z_*+1)\omega(h))$ integers by roughness. For $q>L^B$, one-boundedness gives an error $O(U/q+1+U/\varphi(q))$. If $q\le U^2$, use $q/\varphi(q)\ll\log\log(3q)\ll\log L$ and take $B$ sufficiently large. If $q>U^2$, then $\varphi(q)\ge\sqrt{q/2}>U/\sqrt2$, so the error is $O(1)$, admissible since $U\ge X^{2\zeta}$. The totient bound follows from $\varphi(p^e)^2\ge p^e$ on prime powers except $2$. Truncating the complementary divisor-bounded sequence at $L^B$ via~\eqref{eq:shifted_div} and applying Lemma~\ref{lem:packing}, Pascadi~\cite[Proposition~4.4]{Pascadi2025} (after dyadically subdividing the three modulus groups, since all powers in~\eqref{eq:pascadi_cond} are positive) yields~\eqref{eq:distrib}.

\smallskip
\noindent\emph{Case 2: No medium slot exists.} With $\zeta \ll a_*$, each of the two original primes contains a long slot of length $> X^{1/3+\zeta/2}$ (since six short slots have product $O(X^{12\zeta})$). There are exactly two long slots (three would exceed $X$), neither can be truncated M\"obius, and the short complementary product has scale $M \le X^{24\zeta+o(1)}$ and $L^1$-norm $O(ML^C)$. We prove an absolute modulus average.

For fixed short product $m$ with $(m,d)=1$, set $b_0 \equiv 2\bar m \pmod d$. For smooth compact long profiles $f_0,g_0$ at scales $U,V$, double Poisson summation gives the unrestricted discrepancy
\begin{equation}\label{eq:poisson2}
\frac{UV}{d^2}\sum_{j,k\in\mathbb{Z}}\widehat f_0(jU/d)\widehat g_0(kV/d)
\left(S(b_0 j, k; d) - \frac{c_d(j)c_d(k)}{\varphi(d)}\right).
\end{equation}
Both axes $j=0$ and $k=0$ vanish because $(b_0,d)=1$. Using $|\widehat f_0(t)| \ll (\|f_0\|_1 + \|f_0''\|_1)(1+|t|)^{-2}$, the Weil bound $|S(b_0 j, k; d)| \le \tau(d)\sqrt{d}\,(j,k,d)^{1/2}$, and $\sum_{j \ne 0,\, h\mid j}(1+|j|/H)^{-2} \le 2H/h$, summing over $d \le D_0$ yields $O(M D_0^{3/2} L^{O(1)})$ (and $O(\tau(d)^2/\varphi(d))$ from Ramanujan sums), which has a fixed power saving since $24\zeta + \tfrac{3}{2}(5/8+\xi) < 1$.

To restore roughness without multiplying negative lower-sieve polynomials, set $h_0$ to the even integer $2\lfloor(\log\log X)^2\rfloor$ and put $y_* = z_*^{h_0+1} = X^{o(1)}$. Let $r_z := \ind_{P^-(n)>z_*}$ denote the rough indicator. Form Brun polynomials $L_z$ and $U_z$ satisfying $L_z \le r_z \le U_z$, truncated at $h_0+1$ and $h_0$ prime factors, with $G_z = U_z - L_z \ge 0$ and coefficients bounded by $1$. Their coprime densities satisfy $a_d \ll 1$ and $0 \le g_d \le (\sum_{\ell \le z_*} 1/\ell)^{h_0+1}/(h_0+1)! \ll_A L^{-A}$. Using
\begin{equation}\label{eq:brun_diff}
0 \le U_1 U_2 - r_1 r_2 = (U_1 - r_1)U_2 + r_1(U_2 - r_2) \le G_1 U_2 + U_1 G_2,
\end{equation}
expanding Brun divisors shifts $(U,V,m)$ to $(U/e, V/e', mee')$ with at most a $y_*^2$ loss in~\eqref{eq:poisson2}, while the coprime means from
\[
\sum_{(n,d)=1} f_0(n/U) U_z(n) = U\frac{\varphi(d)}{d}a_d\int f_0 + O\bigl(y_*\tau(d)(\|f_0\|_1 + \|f_0'\|_1)\bigr)
\]
and its $G_z$ analogue give total mean error $L^{O(1)}(MUV\varepsilon_X + M(U+V)y_* + My_*^2) \ll_A XL^{-A}$. Sharp interval indicators are retained via smooth error strips of relative width $L^{-B}$. This completes the proof of~\eqref{eq:distrib}.
\end{proof}

\section{Limiting density and small-prime tail}\label{sec:density}

Let $\normF(s) := F(s)/(2e^\gamma)$ and $\normf(s) := f(s)/(2e^\gamma)$ be the normalized linear-sieve functions, extended by $\normf(s) = 0$ for $s \le 2$. They satisfy
\begin{equation}\label{eq:delay}
\begin{gathered}
\normF(s) = \frac{1}{s}\ (1\le s\le 3),\qquad \normf(s)=0\ (s\le 2),\\
(s\normF(s))' = \normf(s-1)\ (s\ge 3),\qquad (s\normf(s))' = \normF(s-1)\ (s\ge 2).
\end{gathered}
\end{equation}
Here $\normF$ decreases and $\normf$ increases to the same limit, $\normF$ is convex on $s \ge 1$ (as $-\normF' = (\normF - \normf(s-1))/s$ decreases), and $\normf$ is concave on $s \ge 2$ (as $\normf' = (\normF(s-1) - \normf)/s$ decreases), with an upward derivative jump of $1/2$ at $s=2$.

In units of $C_2/L$, the ordinary reference cubic Rosser tree of cap $T = 5/8$ and root $\rho = 1/5$ has density $2\normF(T/\rho)/\rho = 2/T + I_1$, where with $h = T/\rho - 3 = 1/8$,
\begin{equation}\label{eq:I1}
I_1 = \frac{2}{T}\int_0^h \frac{\log(1+t)}{2+t}\,dt \le \frac{h^2}{2T}.
\end{equation}
By the exit-series form of the linear sieve (with uniform prime-harmonic convergence on compact positive parameter ranges, noting that at any fixed exit rank the last exponent is bounded below by a positive multiple of $\rho$~\cite{Iwaniec1980,FriedlanderIwaniec2010}), every third-rank subtree is a negative copy of an ordinary lower cubic sieve of remaining level $J - S_3$ and threshold $t_3$. Hence the rank-three replacement adds
\begin{equation}\label{eq:I3}
I_3 = 2\int_{0 < z < y < u < \rho}
\frac{\normf((T-u-y-z)/z) - \normf((J-u-y-z)/z)}{uy z^2}\,du\,dy\,dz.
\end{equation}
On $z < \delta'$, all finite approximations keep $J - u - y \ge c > 0$, so by the fundamental lemma the residual density difference is $O(V(x^z)e^{-c'/z}) = O(C_2 z/L)$; integrating against the first two harmonic measures $\ll (1 + \log(\rho/z))^2$ and dyadic $z$ gives $O(\delta'(1 + \log^2(\rho/\delta'))C_2/L)$. Thus the compact limit at $z \ge \delta'$ and dominated convergence yield~\eqref{eq:I3} and show that the finite construction and attachment errors in Section~\ref{sec:packing} can be kept at most $2u_0 C_2/L$ for any chosen $u_0 > 0$. There is no separate single-prime top interval above $\mu_0$ because $\rho = 1/5 < 1/4$.

To bound the small-$z$ tail of~\eqref{eq:I3} explicitly, let $\normD := \normF - \normf \ge 0$ and $\delta = 19/1000$. By~\eqref{eq:delay},
\begin{equation}\label{eq:delay_D}
\normD'(s) = \frac{-\normD(s) - \normD(s-1)}{s} \le -\frac{2\normD(s)}{s} \qquad (s \ge 3),
\end{equation}
so $\normD(s) \le \normD(9)(9/s)^2$ for $s \ge 9$. Since $J - u - y \ge (1 - 3\rho)/2 = \rho$, the lower argument in~\eqref{eq:I3} exceeds $(\rho - \delta)/z > 9$ for $z < \delta$. Bounding the numerator by $\normD((J-u-y-z)/z)$ and integrating $z < y < u < \rho$ gives
\begin{equation}\label{eq:I3t}
I_{3,t} \le \frac{81\normD(9)\delta}{(\rho-\delta)^2}\left(\log^2(\rho/\delta) + 2\log(\rho/\delta) + 2\right).
\end{equation}
Letting $I_{3,c}$ be~\eqref{eq:I3} restricted to $z \ge \delta$, the rank-three baseline is $U_0 = 2/T + I_1 + I_{3,c} + I_{3,t}$.

At rank four ($t_4 \ge 3/100$, already compact), the residual subtree is a positive upper sieve, so increasing the cap from $J$ to $M$ subtracts the nonnegative gain
\begin{equation}\label{eq:G1}
G_1 = 2\int_{\mathcal{A}}
\frac{\normF((J - S_4)/t_4) - \normF((M - S_4)/t_4)}{t_1 t_2 t_3 t_4^2}\,dt_1\,dt_2\,dt_3\,dt_4,
\end{equation}
where $\mathcal{A}$ consists of ordered rank-four prefixes passing the third cube and admitting a strict improvement~\eqref{eq:rank4_guard}. Note $(J-S_4)/t_4 > 1$ because $J > S_2 + 3t_3$ implies $J - S_4 > 2t_3 - t_4 \ge t_4$.

\section{Application to $B_b(x)$ and proof of Theorem~\ref{thm:main}}\label{sec:bad_pairs}

Let $K$ be any upper bound for $U_0 + 2u_0$ (rank three) or $U_0 + 2u_0 - G_1$ (ranks three and four). For every row of the certificate table fix $a_*=1/3$ before constructing $\rho_x$, since $b/(1+b)>1/3$. Discard $p<x/L^B$ ($B>4$) from $B_b(x)$ at cost $O(x/L^B)$. For the remaining pairs, $rq\ge x/L^B$ and $r>(rq)^{b/(1+b)}$ give $q\ge r\gg x^{a_*}$, and $p>x^{\rho'}$. Applying Lemma~\ref{lem:first_exit} pointwise to $rq - 2$ and majorizing the pair region by smooth outer weights of polylogarithmic Euler norms, Proposition~\ref{prop:distribution} yields an upper bound controlled by the unrestricted pair mass times $\sum_d \rho_x(d)/\varphi(d)$: replacing $\ind_{(rq,d)=1}$ by $1$ costs $O(L x^{-a_*})$ relative to the pair mass, because $\sum_{d \le 2x,\, \ell \mid d} 1/\varphi(d) \ll L/\ell$ for each $\ell \in \{r,q\}$.

In coordinates $t = \log r/L$ and $y = rq/x$, on $y \ge \sigma > 0$ the inequalities $q^b < r \le q$ and $rq \le x+2$ match $\beta_b := b/(1+b) < t \le 1/2$ up to $O_\sigma(1/L)$, while $y \le \sigma$ is covered by a smooth majorant on $y < 2\sigma$, $t \ge \beta_b - o(1)$, $r \le 2\sqrt{x}$ of mass $O(\sigma x/L)$ via the prime number theorem. Sending the smoothing width and then $\sigma$ to zero gives limiting pair mass $\int_{\beta_b}^{1/2} dt/(t(1-t)) = \log(1/b)$ in units of $x/L$. Combining with~\eqref{eq:sub} gives
\begin{equation}\label{eq:final_sub}
\liminf_{x\to\infty}\frac{A_b(x)}{H(x)}
\ \ge\ \frac{12759}{10000} - K\log(1/b).
\end{equation}

\section{Directed integer certificate}\label{sec:certificate}

The companion script \texttt{verify\_twin\_upper\_sieve\_cert.py} certifies all numerical inequalities in pure integer arithmetic at denominator $Q = 2^{80} = 1208925819614629174706176$ (using \texttt{Fraction}, integer \texttt{isqrt} rounded upward, and floor/ceiling integer divisions, with zero floating-point operations in any bound).

\begin{enumerate}
\item \emph{Logarithms and delay tables.} For $1 \le z \le 2$, $\log z = 2\sum_{j=0}^{47} v^{2j+1}/(2j+1) + R_{48}$ ($v=(z-1)/(z+1)$) with $0 \le R_{48} \le 2v^{97}/(97(1-v^2))$; larger arguments reduce by powers of $2$. The tables for $s\normF$ and $s\normf$ on a $1/4096$ mesh integrate the preceding numerator chord against $ds/s$ using alternating fourth/fifth partial sums of the weights $(j+1)\log(1+1/j)-1$ and $1-j\log(1+1/j)$. Since $s\normF$ is convex and $s\normf$ is concave (each with curvature bounded by $1$ on its nonconstant cells), the chord error $\le 1/(12n^2 j)$ is subtracted in the lower bound for $s\normf$ (for $j \ge 3n$) and added in the upper bound for $s\normF$. Beyond $s = 12$, both functions lie in $[\normf(12), \normF(12)]$. An independent $44$-term quarter-interval Taylor solver (\texttt{reference\_check}) verifies containment of both branches at $s = 4, \dots, 12$ within $10^{-12}$.

\item \emph{Ordered simplex geometry and kinks.} On any grid box intersected with $t_1 \ge t_2 \ge \dots \ge t_d$, each run of $m$ coincident grid intervals forms an ordered simplex of volume fraction $1/m!$, order-statistic centroid $l + j(h-l)/(m+1)$, and coordinate variance $\le (h-l)^2/12$. For an affine form $Y$, $\mathbb{E}Y_+ \le (\sqrt{\operatorname{Var}(Y) + (\mathbb{E}Y)^2} + \mathbb{E}Y)/2$, using Cauchy--Schwarz within each run, independence across runs, and $\sqrt{12} > 17/5$. At rank three, only branches whose strict guards hold throughout the cell enter the lower minorant of $J$; at rank four, any unrefuted parent branch enters the upper majorant of $J$, while a candidate branch~\eqref{eq:rank4_guard} enters the lower gain only when both planes satisfy all guards throughout the cell (using the single highest certified gain per cell, with $\normF(s) = 1/s$ extended to $0 < s \le 1$).

\item \emph{Cell Taylor bounds.} For a proxy affine cap $P$ with residual ratio $s = (P-S_d)/t_d$ and centered deviation $g = P - S_d - s(\mathbf{c})t_d$, the weighted linear Taylor error on $\phi \in \{\normF, \normf\}$ is bounded by
\[
B(\mathbf{l})\frac{|\phi'(s(\mathbf{c}))|\,\sigma_g}{l_d}\sum_i (a_i + \ind_{i=d})\frac{\sigma_i}{l_i},
\]
where $B(\mathbf{t}) = 2\prod_i t_i^{-a_i}$ ($a_i = 1 + \ind_{i=d}$), tightened on differences of same-sign slopes via $d_1\sigma_{g_1-g_2} + |d_1-d_2|\sigma_{g_2}$ and its symmetric counterpart. Quadratic curvature corrections $\le B(\mathbf{l})\sup|\phi''|\,\sigma_g^2/(2l_d^2)$ are charged in the non-convex/non-concave direction, with an extra $B(\mathbf{l})\sigma_g/(4l_d)$ whenever the upper estimate for $\normf$ straddles the $s=2$ kink. The common weight $B$ satisfies $B(\mathbf{c}) \le \mathbb{E}B \le B(\mathbf{c}) + \frac{1}{2}B(\mathbf{l})[(\sum_i a_i\sigma_i/l_i)^2 + \sum_i a_i(\sigma_i/l_i)^2]$.
\end{enumerate}

\subsection{Certified rational bounds}
Taking $u_0 = 1/10^6$, running the verifier evaluates $435001$ third-rank and $700001$ fourth-rank cells, exits with code~$0$, and certifies the following exact numerators at $Q = 2^{80}$:

\begin{table}[ht]
\centering
\small
\begin{tabular}{@{}lcr@{}}
\toprule
Quantity & Direction & Exact numerator at $Q=2^{80}$ \\
\midrule
$I_1$ & $\le$ & \texttt{13943390692550619486377} \\
$I_{3,c}$ & $\le$ & \texttt{66925635371004488689726} \\
$I_{3,t}$ & $\le$ & \texttt{3409294078124169319} \\
$G_1$ & $\ge$ & \texttt{19374509949983140898664} \\
$K_3$ (for $U_0+2u_0$) & $\le$ & \texttt{3949437475976085820663535} \\
$K_4$ (for $U_0+2u_0-G_1$) & $\le$ & \texttt{3930062966026102679764871} \\
\bottomrule
\end{tabular}
\caption{Certified rational density components ($K_4 \approx 3.2508717$).}
\end{table}

\begin{table}[ht]
\centering
\small
\begin{tabular}{@{}llr@{}}
\toprule
$b = a - 1$ & Construction & Lower bound numerator on margin at $Q=2^{80}$ \\
\midrule
$3539/5000$ ($1-1.7078$) & Two-factor shortcut & \texttt{377883443722357593132} \\
$6767/10000$ ($1-1.6767$) & Rank three ($K_3$) & \texttt{105553009773095358394} \\
$17/25$ ($1-1.68$) & Ranks three and four ($K_4$) & \texttt{26790620021371983365686} \\
$169/250$ ($1-1.676$) & Ranks three and four ($K_4$) & \texttt{3604340579464536079328} \\
$3377/5000$ ($1-1.6754$) & Ranks three and four ($K_4$) & \texttt{114570068060210771664} \\
$67539/100000$ ($1-1.67539$) & Ranks three and four ($K_4$) & \texttt{56380964408304863785} \\
\bottomrule
\end{tabular}
\caption{Certified positive lower margins for $\frac{12759}{10000} - K\log(1/b)$ in~\eqref{eq:final_sub}.}
\end{table}

All six margins are strictly positive; in particular, the final row gives~\eqref{eq:main_margin} and completes the proof of Theorem~\ref{thm:main}.

\section*{Reproducibility and Acknowledgments}
The standalone Python verifier \texttt{verify\_twin\_upper\_sieve\_cert.py} runs with the Python standard library alone (no external packages required) and verifies every bound in Tables~1 and~2 in approximately three minutes. The analytic derivations and certificate implementation were developed in an AI-assisted human-in-the-loop research workflow, with independent line-by-line mathematical auditing and full re-execution of the integer verifier.

\begin{thebibliography}{99}

\bibitem{BakerHarman1996}
R.~C.~Baker and G.~Harman,
\emph{The Brun--Titchmarsh theorem on average},
in \emph{Analytic Number Theory, Vol.~1 (Allerton Park, IL, 1995)},
Progr. Math.~\textbf{138}, Birkh\"auser, Boston, 1996, pp.~39--103.

\bibitem{Chen1973}
J.~R.~Chen,
\emph{On the representation of a larger even integer as the sum of a prime and the product of at most two primes},
Sci. Sinica~\textbf{16} (1973), 157--176.

\bibitem{DeshouillersIwaniec1982}
J.-M.~Deshouillers and H.~Iwaniec,
\emph{Kloosterman sums and Fourier coefficients of cusp forms},
Invent. Math.~\textbf{70} (1982), no.~2, 219--288.

\bibitem{FriedlanderIwaniec2010}
J.~Friedlander and H.~Iwaniec,
\emph{Opera de Cribro},
Amer. Math. Soc. Colloq. Publ.~\textbf{57}, Amer. Math. Soc., Providence, RI, 2010.

\bibitem{Iwaniec1980}
H.~Iwaniec,
\emph{Rosser's sieve},
Acta Arith.~\textbf{36} (1980), 171--202.

\bibitem{Iwaniec1980b}
H.~Iwaniec,
\emph{A new form of the error term in the linear sieve},
Acta Arith.~\textbf{37} (1980), 307--320.

\bibitem{RunboLi2026}
R.~Li,
\emph{On Chen's theorem, Goldbach's conjecture and almost prime twins~II},
Math. Reports~\textbf{28} (2026), no.~1--2, 39--61; arXiv:2405.05727v4.

\bibitem{LiLiu2026}
J.~Li and J.~Liu,
\emph{Theorem $(1+1.9)$ on the Goldbach conjecture},
preprint, arXiv:2606.05224v2 (2026).

\bibitem{Chan2026}
S.~Chan,
\emph{Theorem $(1+1.75)$ on the Goldbach conjecture},
preprint, aiXiv:261111.000001v1.0 (2026),
\url{https://aixiv.science/pdf/aixiv.261111.000001v1.0.pdf}.

\bibitem{Lichtman2023}
J.~D.~Lichtman,
\emph{Primes in arithmetic progressions to large moduli, and Goldbach beyond the square-root barrier},
preprint, arXiv:2309.08522v1 (2023).

\bibitem{Pascadi2025}
A.~Pascadi,
\emph{On the exponents of distribution of primes and smooth numbers},
preprint, arXiv:2505.00653v1 (2025).

\end{thebibliography}

\end{document}
